Peephole optimization is a key class of compiler optimizations that
targets small instruction sequences, typically within a basic block,
to improve performance and reduce code size. These optimizations are
usually implemented as pattern-matching rules that identify specific
instruction sequences and replace them with more efficient alternatives.
Ensuring the correctness of these transformations is critical, as any
incorrect optimization can cause erroneous program behavior.

In \llvm, a widely used compiler infrastructure, peephole optimizations
are mainly implemented in the \instcombine pass. As of this writing,
\instcombine contains over 50{,}000 lines of code—double its size from
ten years ago. Its growing complexity makes maintenance and correctness
verification increasingly difficult.
Yang \etal found \instcombine to be among the most bug-prone components
in \llvm~\cite{yang2011finding}.
This underscores the need for more rigorous correctness guarantees
for peephole optimizations in \llvm.

\alive~\cite{lopes2021alive2} is a formal verification framework based on \smt solvers,
specifically designed to verify \llvm optimizations. However, \alive
has several limitations, including timeouts for complex optimizations
and lack of support for arbitrary-width bit vectors.

